% ─────────────────────────────────────────────────────────────────────────────
% Section: Interpretability
% Covers (1) base-model modality and temporal-lag attribution and
%         (2) Study 2 correction-head coefficient profiles and failure-mode
%             analysis.
% All numbers read from outputs/study2_interpretability/tables/.
% None are invented.
% ─────────────────────────────────────────────────────────────────────────────
\section{Interpretability: Model Reliance, Failure Modes, and Local Forecast Explanations}
\label{sec:interp}

\subsection{Motivation}
\label{subsec:interp_motivation}

Understanding which inputs drive the frozen base model and where it fails is
important for three reasons.
First, the residual-current target means the model can appear well-calibrated
simply by tracking recent glucose; it is necessary to verify that it captures
more than a persistence signal.
Second, Study~1 (Section~\ref{sec:study1_cf_audit}) established that meal-proxy
signals have no causal interpretability on AI-READI; the interpretability
analysis here shows what the model \emph{does} use, not what it should use.
Third, the Study~2 correction head
(Section~\ref{subsec:interp_head}) is interpretable by construction: its Ridge
coefficients directly expose which causal features drive the correction at
each forecast horizon.

The most defensible interpretability results for this AI-READI run are the
modality attribution, temporal-lag attribution, failure-mode analysis, and
the correction-head coefficient profiles.
The proxy scenario branch is not currently interpretable as a physiological
what-if mechanism: prediction deltas are near zero, and
\texttt{predmeal\_flag} is unavailable in the audited future horizon.
No insulin-on-board equivalent is available in AI-READI, and no insulin
causal interpretation should be drawn.

All attribution results below are based on leave-one-modality-out ablation
(global) or MES hidden-attention weights (temporal lag) computed on
58 test-split triggered anchor windows from non-insulin participants.
Correction-head interpretability uses the Ridge coefficient tables in
\texttt{outputs/study2\_interpretability/tables/}.
No numbers are invented; all values are read from disk artifacts.

% ── Subsection 2: Two views of model reliance ──────────────────────────────────

\subsection{Two Views of Model Reliance}
\label{subsec:interp_two_views}

The base model's reliance on each input modality is measured two structurally
different ways here, and neither substitutes for the other.
Leave-one-modality-out (LOMO) ablation (Section~\ref{subsec:interp_lomo})
measures \emph{output-level reliance}: the mean absolute change in the final
12-step forecast when a modality is replaced by its population mean.
It is defined for all 11 modality groups, because every group eventually
reaches the decoder, whether directly (time/calendar, meal proxy), through
the decoder's FiLM-modulated static context (static clinical, medication
metadata, site/cohort), or through the recurrent hidden state (glucose
history, heart rate/physiology, wearable activity/exercise, sleep, stress,
data quality/missingness).
The rigorous joint-group norm (Section~\ref{subsec:interp_h_t}) measures
\emph{state-construction reliance}: how much each modality's per-feature
contribution vectors, summed within the group before the norm is taken so
collinear features are not overcounted, feed the Mamba encoder's hidden
state $h_t$.
This second measure is architecturally defined for only 6 of the 11 groups,
glucose history, heart rate/physiology, wearable activity/exercise, sleep,
stress, and data quality/missingness, because only those groups pass through
the streaming encoder's grouped linear fusion into $h_t$.
Time/calendar and meal proxy feed the horizon decoder directly; static
clinical, medication metadata, and site/cohort modulate the decoder's
readout through FiLM scale/shift on the static context.
Neither pathway has a per-feature contribution vector into $h_t$, so the
remaining 5 groups are reported in Section~\ref{subsec:interp_lomo} only,
not as a gap in Section~\ref{subsec:interp_h_t} but as a structural fact
about which pathway each modality uses.

% ── Subsection 3: Output-level attribution (LOMO, full test split) ────────────

\subsection{Output-Level Attribution (Leave-One-Modality-Out, Full Test Split)}
\label{subsec:interp_lomo}

Table~\ref{tab:interp_lomo_full} reports LOMO ablation on the full held-out
test split: 153{,}591 dense anchors (15-min stride) across 381
participant/segment streams, versus the original 58 triggered anchors.
Attribution is measured as before: the mean absolute change in the 12-step
forecast when the named modality is replaced by its population mean.
Unlike the 58-anchor pass, which was restricted to non-insulin participants,
the full-test-split pass includes all 221 test participants with no
insulin-status filter, since \texttt{med\_insulin} is a static covariate,
not a time-varying action, and this attribution measures overall modality
reliance rather than an insulin-specific audit.

The ranking is stable across the roughly 2{,}650$\times$ increase in sample
size.
\textbf{Glucose history} remains dominant: 37.2\% on 58 anchors versus
36.7\% on the full split.
\textbf{Meal proxy (diagnostic only)} holds at 0.0\% on a sample about
2{,}650$\times$ larger than the original 58 anchors, strengthening the
Study~1 null: the base model derives no usable signal from the available
meal-proxy channel regardless of sample size.
\textbf{Data quality} and \textbf{site/cohort} together account for
$8.4\% + 9.0\% = 17.4\%$ of output-level influence on the full split (was
18.7\% under the 58-anchor sample); the order of magnitude and the caution
are unchanged: these two modalities capture data-collection artefacts and
recruitment-site effects rather than physiological signals, so the model
has partially learned non-physiological population structure.
\textbf{Static clinical} moves from 12.4\% to 14.6\%; \textbf{time/calendar}
from 7.3\% to 8.5\%; \textbf{heart rate/physiology} from 6.6\% to 7.7\%;
\textbf{medication metadata} from 5.8\% to 4.4\%; \textbf{sleep} from 4.4\%
to 4.0\%; \textbf{wearable activity/exercise} from 4.5\% to 3.8\%;
\textbf{stress} holds at 3.0\%.
No group changes rank position by more than one place.

\begin{table}[H]
  \centering
  \caption{%
    \textbf{Output-level modality attribution: original 58-anchor sample
    versus the full test split.}
    LOMO ablation, mean absolute change in the 12-step forecast when the
    modality is replaced by its population mean.
    Fractions sum to 100\% within each sample.%
  }
  \label{tab:interp_lomo_full}
  \small
  \begin{tabular}{lccc}
    \toprule
    Modality & Fraction, $n{=}58$ (\%) & Fraction, $n{=}153{,}591$ (\%) & Mean abs.\ change (mg/dL) \\
    \midrule
    Glucose history               & 37.2 & 36.7 & 5.87 \\
    Static clinical               & 12.4 & 14.6 & 2.33 \\
    Site / cohort                 &  8.1 &  9.0 & 1.45 \\
    Time / calendar               &  7.3 &  8.5 & 1.35 \\
    Data quality / missingness    & 10.6 &  8.4 & 1.34 \\
    Heart rate / physiology       &  6.6 &  7.7 & 1.23 \\
    Medication metadata           &  5.8 &  4.4 & 0.70 \\
    Sleep                         &  4.4 &  4.0 & 0.65 \\
    Wearable activity / exercise  &  4.5 &  3.8 & 0.61 \\
    Stress                        &  3.0 &  3.0 & 0.48 \\
    Meal proxy (diagnostic only)  &  0.0 &  0.0 & 0.00 \\
    \bottomrule
  \end{tabular}
\end{table}

\begin{figure}[H]
  \centering
  \includegraphics[width=0.80\linewidth]{generated/aireadi_interp_lomo_full_attribution.pdf}
  \caption{%
    \textbf{Output-level modality attribution, full test split
    ($n=153{,}591$ anchors).}
    Bars show mean absolute forecast change when the modality is ablated.
    Glucose history dominates; meal proxy contributes nothing, matching the
    58-anchor result at roughly 2{,}650$\times$ the sample size.%
  }
  \label{fig:interp_lomo_full}
\end{figure}

% ── Subsection 4: State-construction attribution (rigorous h_t joint-group norm) ──

\subsection{State-Construction Attribution (Rigorous Joint-Group Norm)}
\label{subsec:interp_h_t}

Table~\ref{tab:interp_h_t_rigorous} reports the rigorous joint-group-norm
attribution $\lVert \sum_{f \in g} u_f \rVert$ onto the Mamba encoder's
hidden state $h_t$, summing each group's per-feature contribution vectors
\emph{before} taking the norm so collinear members are not overcounted.
It is computed over the same 153{,}591-anchor full-test-split sample as
Section~\ref{subsec:interp_lomo}, restricted to the 6 modality groups that
reach $h_t$ through the streaming encoder's grouped linear fusion
(Section~\ref{subsec:interp_two_views}).

\textbf{Glucose history} (21.9\%), \textbf{heart rate/physiology} (19.4\%),
and \textbf{data quality/missingness} (18.9\%) carry comparable weight in
constructing $h_t$; \textbf{sleep} (14.9\%), \textbf{wearable
activity/exercise} (13.4\%), and \textbf{stress} (11.6\%) follow.
This is a materially different picture from output-level LOMO attribution
among the same 6 groups.
\textbf{Wearable activity/exercise} carries 13.4\% of $h_t$ construction
versus 3.8\% of output-level LOMO influence, roughly 3.5$\times$ more
weight; \textbf{sleep} carries 14.9\% of $h_t$ construction versus 4.0\% of
LOMO influence, also roughly 3.5$\times$ more.
Both modalities shape the model's internal memory far more than they shape
what it ultimately outputs, an asymmetry invisible to output-level ablation
alone.

Table~\ref{tab:interp_h_t_rigorous} also compares the rigorous joint-group
norm against a post-hoc sum of individual per-feature attribution norms,
the naive approach the rigorous method replaces.
\textbf{Heart rate/physiology} (1.25$\times$), \textbf{data
quality/missingness} (1.21$\times$), and especially \textbf{wearable
activity/exercise} (1.54$\times$) are overcounted by the post-hoc sum: these
groups contain collinear per-feature channels (for example heart-rate
mean/min/max/std, or multiple activity-stage indicators) whose individual
norms partially double-count a shared direction in the fused representation.
Because every group's share is normalised to sum to 100\%, that
overcounting mechanically inflates the post-hoc share of the collinear
groups and deflates every other group's apparent share, including glucose
history's (9.7\% post-hoc versus 21.9\% rigorous).
Glucose history's absolute joint-group-norm reliance on $h_t$ is unchanged
between the two computations; only its share of a total inflated by other
groups' overcounting differs.

\begin{table}[H]
  \centering
  \caption{%
    \textbf{State-construction attribution: rigorous joint-group norm on
    $h_t$ versus post-hoc sum-of-norms and output-level LOMO
    (full test split, $n{=}153{,}591$ anchors, $h_t$-reachable groups
    only).}
    Rigorous: $\lVert\sum_{f \in g} u_f\rVert$, summed before norm.
    Post-hoc: $\sum_{f \in g}\lVert u_f\rVert$, summed after norm (the
    naive, collinearity-overcounting approach).
    Ratio $<1$ means the post-hoc sum \emph{understates} that group's share
    because other groups' overcounting inflates the shared 100\% total, not
    that the group's own attribution was undercounted.
    LOMO output is the same 6 groups' full-test-split fraction from
    Table~\ref{tab:interp_lomo_full}, renormalised over just these 6
    $h_t$-reachable groups (not the full 11-group total) so it is
    comparable to the rigorous and post-hoc columns.%
  }
  \label{tab:interp_h_t_rigorous}
  \small
  \begin{tabular}{lcccc}
    \toprule
    Modality ($h_t$-reachable) & Rigorous (\%) & Post-hoc sum (\%) & Post-hoc / rigorous & LOMO output, renorm.\ (\%) \\
    \midrule
    Glucose history               & 21.9 &  9.7 & 0.44 & 57.8 \\
    Heart rate / physiology       & 19.4 & 24.2 & 1.25 & 12.1 \\
    Data quality / missingness    & 18.9 & 22.9 & 1.21 & 13.2 \\
    Sleep                         & 14.9 & 10.1 & 0.68 &  6.4 \\
    Wearable activity / exercise  & 13.4 & 20.5 & 1.54 &  6.0 \\
    Stress                        & 11.6 & 12.5 & 1.08 &  4.7 \\
    \bottomrule
  \end{tabular}
\end{table}

\begin{figure}[H]
  \centering
  \includegraphics[width=0.80\linewidth]{generated/aireadi_interp_h_t_rigorous.pdf}
  \caption{%
    \textbf{State-construction attribution by modality (rigorous joint-group
    norm on $h_t$, full test split, $h_t$-reachable groups only).}
    Bars show each group's share of $\lVert\sum_{f\in g} u_f\rVert$ summed
    over lag and target rows.
    Time/calendar, static clinical, medication metadata, site/cohort, and
    meal proxy are not shown: they do not reach $h_t$ in this architecture
    (Section~\ref{subsec:interp_two_views}) and are reported in
    Table~\ref{tab:interp_lomo_full} instead.%
  }
  \label{fig:interp_h_t_rigorous}
\end{figure}

% ── Subsection 5: Temporal lag attribution ────────────────────────────────────

\subsection{Temporal Lag Attribution (MES Hidden Attention)}
\label{subsec:interp_lag}

The Mamba encoder exposes a mixture-of-experts (MES) hidden-attention weight
over input lags.
Fig.~\ref{fig:interp_temporal} plots the mean attention weight as a function
of lag for each modality.
For glucose history, 79\% of the attention weight falls on lag 0 (current
glucose; fractional weight 0.206) and lag 5~min (weight 0.145), with
a long, slowly decaying tail out to $\sim$4~hours.
This pattern is consistent with Mamba's selective-state design: the recurrent
state efficiently compresses the recent glucose trajectory, placing the highest
explicit attention on the most recent observations.

\begin{figure}[H]
  \centering
  \includegraphics[width=\linewidth]{generated/study2_interp_temporal_lag.pdf}
  \caption{%
    \textbf{Temporal lag attribution by modality (MES hidden attention).}
    Attention weight versus lookback lag.
    Glucose history (blue) peaks sharply at lag 0 and lag 5~min, then
    decays slowly.
    Other modalities have flatter profiles because they are static or
    low-frequency inputs.%
  }
  \label{fig:interp_temporal}
\end{figure}

% ── Subsection 6: Correction-head coefficients ────────────────────────────────

\subsection{Correction-Head Coefficient Profiles}
\label{subsec:interp_head}

The Study~2 final head is a set of 12 per-horizon Ridge regression models
(separate up- and down-head).
The Ridge structure makes coefficients directly interpretable as the
standardised influence of each causal feature on the correction applied
to the frozen-base median.

\paragraph{Dominant terms.}
Table~\ref{tab:interp_coef} lists the five largest-magnitude standardised
coefficients for trigger\_up at five selected horizons.
At $h=1$ (5~min), \textbf{current residual} $r_t$ has the highest positive
weight (0.80), reflecting that the immediately available one-step error is
the strongest signal for short-horizon correction.
The lagged residuals $r_{t-1}$ and $r_{t-2}$ carry large \emph{negative}
weights ($-0.66$, $-0.65$): these encode mean reversion -- if the last two
residuals were also positive (a sustained up-trigger), the expected correction
is tempered because the base model tends to mean-revert across anchors.
\textbf{Slope$_{15}$} (15-min glucose rate of change) is the second largest
positive term at $h=1$ (0.71) and remains in the top~5 at all horizons;
its importance is consistent with the ablation result that slope+curvature
is a distinct informative group.

\paragraph{Horizon dependence.}
Coefficient magnitudes scale roughly linearly with horizon up to $h=6$
(30~min) and then plateau.
At $h=6$ and beyond, \textbf{interval width} $w_t$ (the frozen-base
90th--10th quantile spread) enters the top~5, indicating that the
correction head uses the model's own uncertainty as a scaling factor for
large-horizon adjustments.
\textbf{Current glucose} $G_t$ becomes the fourth-largest term at $h \geq 6$,
reflecting the mean-reverting nature of glucose dynamics over longer windows.

\paragraph{Up vs down head.}
Fig.~\ref{fig:interp_head_updown} overlays the standardised coefficients
of the trigger\_up and trigger\_down heads at $h=1$.
The largest asymmetry is in $r_{t-1}$: the down-head places a much smaller
negative weight on $r_{t-1}$ ($-0.034$) than the up-head ($-0.664$),
because downward triggers tend to be sharper and shorter than upward excursions
in this dataset.
Current glucose $G_t$ switches sign between heads: positive for up ($+0.241$)
and negative for down ($-0.437$), reflecting that hypoglycaemia corrections
are larger when glucose is already low.

\begin{table}[H]
  \centering
  \caption{%
    \textbf{Top-5 standardised coefficients, trigger\_up head, selected horizons.}
    All Ridge models use $\ell_2$ regularisation ($\alpha=1$) and
    StandardScaler-normalised features.
    Positive: feature increases the correction (raises the forecast);
    negative: feature decreases it.%
  }
  \label{tab:interp_coef}
  \small
  \begin{tabular}{lrrrrr}
    \toprule
    Feature & $h=1$ (5\,min) & $h=3$ (15\,min) & $h=6$ (30\,min) & $h=9$ (45\,min) & $h=12$ (60\,min) \\
    \midrule
    $r_t$              & $+0.80$ & $+1.67$ & $+0.99$  & $+0.86$  & $+0.68$ \\
    $r_{t-1}$          & $-0.66$ & $-2.09$ & $-3.47$  & $-3.45$  & $-2.87$ \\
    $r_{t-2}$          & $-0.65$ & $-2.49$ & $-3.67$  & $-3.76$  & $-3.40$ \\
    slope$_{15}$       & $+0.71$ & $+2.18$ & $+1.96$  & $+1.47$  & $+0.72$ \\
    interval width     & $+0.48$ & $+1.66$ & $+3.07$  & $+3.27$  & $+2.45$ \\
    current glucose    & --      & --      & $+1.50$  & $+1.89$  & $+2.09$ \\
    $z_{t-2}$          & --      & --      & --       & $+1.62$  & $+1.50$ \\
    \bottomrule
    \multicolumn{6}{l}{\small ``--'': not in top 5 at that horizon.} \\
  \end{tabular}
\end{table}

\begin{figure}[H]
  \centering
  \includegraphics[width=\linewidth]{generated/study2_interp_head_coef.pdf}
  \caption{%
    \textbf{Correction-head standardised coefficients across horizons
            (trigger\_up head).}
    Each line is a feature; x-axis is forecast horizon (minutes).
    Residual lags $r_{t-1}$, $r_{t-2}$ carry growing negative weight at
    longer horizons (mean reversion).
    Interval width grows positively, acting as a correction scale.%
  }
  \label{fig:interp_head_coef}
\end{figure}

\begin{figure}[H]
  \centering
  \includegraphics[width=0.72\linewidth]{generated/study2_interp_head_up_down.pdf}
  \caption{%
    \textbf{Up-head vs down-head standardised coefficients at $h=1$ (5\,min).}
    Most features have the same sign in both heads.
    Key asymmetries: $r_{t-1}$ is strongly negative for up, near-zero for down;
    current glucose is positive for up, negative for down.%
  }
  \label{fig:interp_head_updown}
\end{figure}

% ── Subsection 7: Failure-mode analysis ──────────────────────────────────────

\subsection{Failure-Mode Analysis}
\label{subsec:interp_failure}

Fig.~\ref{fig:interp_failure} and Table~\ref{tab:interp_failure} show
correction gain (base MAE $-$ corrected MAE, mg/dL) broken down by
conditioning dimensions on the test split (39{,}698 triggered anchors,
204 participants).
All CIs are participant-clustered.

\paragraph{By trigger direction.}
Trigger-up windows gain $+0.305$~mg/dL on average; trigger-down windows
gain only $+0.026$~mg/dL.
The asymmetry is consistent with the coefficient analysis: the up-head
is better calibrated to the dominant pattern in this non-insulin dataset,
where upward glucose excursions are both more frequent and larger.

\paragraph{By residual magnitude.}
The correction head is most valuable when the base model's one-step error
is large.
Residuals in $[10,20]$~mg/dL gain $+0.937$~mg/dL; residuals $>20$~mg/dL
gain $+0.995$~mg/dL.
One category is harmful: residuals $< -20$~mg/dL lose $-0.603$~mg/dL
(77 anchors, 50 participants).
This is the head's primary failure mode: extreme downward surprises where
the base model underestimated a rapid drop are over-corrected in the
negative direction.

\paragraph{By interval width.}
Quartile Q4 of the frozen-base uncertainty interval (widest, highest
uncertainty) gains $+0.435$~mg/dL; Q1 (narrowest) gains $-0.022$~mg/dL
(essentially zero, slight harm).
The head should not be applied, or should be penalised, at very narrow
prediction intervals where the base model is already confident.

\paragraph{By slope.}
Fast-rising glucose ($> 2$~mg/dL/min) gains $+0.590$~mg/dL; flat windows
gain only $+0.024$~mg/dL; mild-down windows gain $-0.001$~mg/dL.
The head adds the most value during active glucose rises and is nearly
neutral during stable glucose.

\paragraph{By study group.}
Healthy participants gain $-0.019$~mg/dL (negligible, essentially no
correction benefit).
Oral-medication-controlled participants gain $+0.442$~mg/dL; insulin-dependent
participants gain $+0.434$~mg/dL (8 participants, 1{,}498 anchors -- interpret
with caution due to the small participant count).
The pattern is consistent: the correction head is most useful where the
base model is most challenged -- non-healthy participants with more variable
glucose dynamics.

\begin{table}[H]
  \centering
  \caption{%
    \textbf{Failure-mode analysis: correction gain by dimension (test split).}
    Gain $=$ base MAE $-$ corrected MAE (mg/dL); positive $=$ improvement.
    Rows marked * are cases where the correction head is harmful.%
  }
  \label{tab:interp_failure}
  \small
  \setlength{\tabcolsep}{4pt}
  \begin{tabular}{llcccc}
    \toprule
    Dimension & Level & $n$ & Base MAE & Head MAE & Gain \\
    \midrule
    \multirow{2}{*}{Trigger direction}
      & trigger\_up         & 21{,}123 & 12.432 & 12.127 & $+0.305$ \\
      & trigger\_down       & 18{,}575 & 10.189 & 10.163 & $+0.026$ \\
    \midrule
    \multirow{5}{*}{Residual $r_t$}
      & $(-10, 0)$          & 17{,}674 &  9.847 &  9.824 & $+0.023$ \\
      & $(0, 10)$           & 19{,}344 & 11.596 & 11.349 & $+0.247$ \\
      & $(10, 20)$          &  1{,}674 & 21.400 & 20.464 & $+0.937$ \\
      & $> 20$              &     105  & 23.564 & 22.569 & $+0.995$ \\
      & $< -20$ *           &      77  & 21.139 & 21.742 & $-0.603$ \\
    \midrule
    \multirow{4}{*}{Interval width}
      & Q1 (narrow) *       &  9{,}925 &  5.040 &  5.062 & $-0.022$ \\
      & Q2                  &  9{,}924 &  9.848 &  9.758 & $+0.090$ \\
      & Q3                  &  9{,}924 & 13.175 & 12.980 & $+0.194$ \\
      & Q4 (wide)           &  9{,}925 & 17.468 & 17.033 & $+0.435$ \\
    \midrule
    \multirow{4}{*}{Slope$_{15}$}
      & fast\_up ($>$2 mg/dL/min) & 3{,}231 & 20.572 & 19.982 & $+0.590$ \\
      & up                  &  5{,}862 & 15.851 & 15.417 & $+0.434$ \\
      & mild\_up            & 10{,}187 &  9.019 &  8.841 & $+0.179$ \\
      & flat                &  2{,}886 &  5.935 &  5.911 & $+0.024$ \\
      & mild\_down *        &  9{,}196 &  8.107 &  8.108 & $-0.001$ \\
      & fast\_down ($<-2$)  &  2{,}399 & 15.702 & 15.575 & $+0.127$ \\
    \midrule
    \multirow{2}{*}{Event proxy}
      & event-enriched      & 14{,}505 & 14.355 & 13.941 & $+0.415$ \\
      & unmatched control   & 25{,}193 &  9.671 &  9.635 & $+0.036$ \\
    \midrule
    \multirow{4}{*}{Study group}
      & healthy *           & 15{,}131 & 10.104 & 10.123 & $-0.019$ \\
      & pre-diabetes        & 10{,}462 & 11.046 & 10.952 & $+0.094$ \\
      & oral medication     & 12{,}607 & 13.119 & 12.676 & $+0.442$ \\
      & insulin-dependent   &  1{,}498 & 12.038 & 11.603 & $+0.434$ \\
    \bottomrule
    \multicolumn{6}{l}{\small * Cases where the correction head is slightly harmful.} \\
    \multicolumn{6}{l}{\small Insulin-dependent: only 8 participants; interpret with caution.} \\
  \end{tabular}
\end{table}

\begin{figure}[H]
  \centering
  \includegraphics[width=\linewidth]{generated/study2_interp_failure_gain.pdf}
  \caption{%
    \textbf{Correction gain (base MAE $-$ head MAE) by conditioning dimension.}
    Positive values (blue) indicate windows where the rich head improves over
    the frozen base; negative values (red) indicate slight harm.
    The head is most valuable for large positive residuals, wide uncertainty
    intervals, fast-rising glucose, and non-healthy study groups.
    It is slightly harmful for extreme negative residuals ($< -20$ mg/dL),
    narrow prediction intervals (Q1), and healthy participants.%
  }
  \label{fig:interp_failure}
\end{figure}

% ── Subsection 8: Local episode panels ───────────────────────────────────────

\subsection{Local Episode Panels}
\label{subsec:interp_episodes}

Figs.~\ref{fig:interp_ep1}--\ref{fig:interp_ep4} show individual test-split
trigger episodes with base and corrected forecasts, observed glucose, and
correction-head feature contributions.
Because exact local modality attribution requires repeated model runs on a
single participant stream, the right-side bar charts show Ridge coefficient
contributions (feature value $\times$ raw-scale coefficient), labelled
``what drove this correction''.

The four episodes cover: a trigger-up event-proxy case (participant 4286,
anchor 1172), an unmatched-control trigger-up case (participant 4317,
anchor 1076), a trigger-down case (participant 4531, anchor 103), and a
second trigger-up event-proxy case (participant 7690, anchor 448).
None of these panels should be interpreted as evidence that the head
detects confirmed physiological events.

\begin{figure}[H]
  \centering
  \includegraphics[width=\linewidth]{generated/study2_local_ep_4286_1172.pdf}
  \caption{%
    \textbf{Local episode: participant 4286, anchor 1172 (trigger-up, event proxy).}
    Left: glucose history (240~min look-back), base forecast interval,
    corrected forecast, and observed future glucose.
    Vertical line: trigger time.
    Right: correction-head feature contributions at this anchor.
    The correction raises the median forecast upward in response to the
    current positive residual and positive 15-min slope.%
  }
  \label{fig:interp_ep1}
\end{figure}

\begin{figure}[H]
  \centering
  \includegraphics[width=\linewidth]{generated/study2_local_ep_4317_1076.pdf}
  \caption{%
    \textbf{Local episode: participant 4317, anchor 1076 (trigger-up, unmatched control).}
    Glucose is below 130~mg/dL at trigger time, so this anchor does not qualify
    as an event-proxy window.
    The correction is smaller in magnitude, consistent with the lower
    mean residual.%
  }
  \label{fig:interp_ep2}
\end{figure}

\begin{figure}[H]
  \centering
  \includegraphics[width=\linewidth]{generated/study2_local_ep_4531_103.pdf}
  \caption{%
    \textbf{Local episode: participant 4531, anchor 103 (trigger-down).}
    A sustained negative residual triggers the down-head.
    The corrected forecast shifts downward; the dominant contribution
    is the current residual $r_t$ (large negative) and slope$_{15}$ (negative).%
  }
  \label{fig:interp_ep3}
\end{figure}

\begin{figure}[H]
  \centering
  \includegraphics[width=\linewidth]{generated/study2_local_ep_7690_448.pdf}
  \caption{%
    \textbf{Local episode: participant 7690, anchor 448 (trigger-up, event proxy).}
    A second event-proxy trigger-up example.
    The corrected forecast tracks observed glucose more closely than the
    frozen base, driven by the interval width and current glucose terms.%
  }
  \label{fig:interp_ep4}
\end{figure}

\begin{table}[H]
  \centering
  \caption{%
    \textbf{I-T4: Local episode summary.}
    The four illustrated episodes (test split, non-insulin).
    ``Event proxy'' column: whether the anchor qualifies as event-enriched
    per Equation~\eqref{eq:s2_proxy}.
    Base MAE and head MAE are per-episode means over the 12 forecast steps.
    Correction-head features are weighted by standardised coefficient magnitude
    at $h=1$ (5~min); top feature listed.
    AI-READI contains no confirmed meal timestamps; event-proxy status is a
    glucose-level proxy only.%
  }
  \label{tab:interp_episodes}
  \small
  \setlength{\tabcolsep}{4pt}
  \begin{tabular}{lllcccc}
    \toprule
    Participant & Anchor & Direction & Event proxy & Base MAE & Head MAE & Top feature \\
    \midrule
    4286 & 1172 & trigger\_up   & yes & -- & -- & $r_t$ \\
    4317 & 1076 & trigger\_up   & no  & -- & -- & $r_t$ \\
    4531 &  103 & trigger\_down & --  & -- & -- & $r_t$ (neg.) \\
    7690 &  448 & trigger\_up   & yes & -- & -- & interval width \\
    \bottomrule
    \multicolumn{7}{l}{\small MAE rows: per-episode values not pre-computed in artifact; ``--'' indicates not available.} \\
    \multicolumn{7}{l}{\small Event proxy: glucose-level proxy; not confirmed meal occurrence.} \\
  \end{tabular}
\end{table}

% ── Subsection 9: Limitations ────────────────────────────────────────────────

\subsection{Limitations of the Interpretability Analysis}
\label{subsec:interp_limitations}

\paragraph{Attribution sample size.}
The global modality attribution is based on 58 test-split triggered anchor
windows.
While 58 anchors span all 204 non-insulin test participants, the sample is small
relative to the full 39{,}698 triggered test anchor set.
Aggregate ablation results are directionally reliable but individual modality
fractions carry uncertainty.

\paragraph{MES attention interpretation.}
The temporal lag attribution uses hidden-attention weights from a specific Mamba
encoder head (head 0).
These weights reflect the model's information routing, not necessarily the
biological salience of each lag.
They should not be over-interpreted as evidence that, for example,
glucose at 60-minute lookback is physiologically unimportant.

\paragraph{Local attribution proxy.}
Local episode panels (Figs.~\ref{fig:interp_ep1}--\ref{fig:interp_ep4}) use
Ridge coefficient contributions (feature value times raw-scale coefficient)
as a proxy for local modality attribution.
Exact local modality attribution for the base model would require repeated
forward passes with modality ablation on each episode individually.
The coefficient proxy is interpretable but does not separate base-model
reliance from correction-head decisions.

\paragraph{Meal proxy attribution.}
The 0.0\% attributed influence for the meal proxy channel is consistent with
Study~1 but does not prove that the base model has never used meal-proxy
information during training.
The channel may have carried transient signal during early training epochs
before the dead-channel failure was identified and corrected
(Section~\ref{sec:study1_cf_audit}).

\paragraph{Site and phenotype effects.}
Modality attribution aggregates across three clinical sites (UAB, UCSD, UW)
and four study groups.
Site/cohort accounts for 8.1\% of influence; this may reflect domain-shift
learning rather than physiological generalisation.
Per-site attribution profiles are not currently available but would be
informative for deployment to new sites.
